% SECCIÓN 3: OBJETIVOS
\section{Objetivos}

% SUB-SECCIÓN
\subsection{Objetivo general}

	Mejorar la continuidad del acceso a Internet para las actividades académicas de estudiantes y docentes en los ambientes del pabellón de la EPIIS de la UNSAAC, a partir de un diagnóstico verificado que permita definir una solución con indicadores medibles.

% SUB-SECCIÓN
\subsection{Objetivos específicos}

	El primer objetivo específico es documentar, con registros fechados por ambiente, la frecuencia y la duración de la falta de señal, los cortes, la lentitud y la necesidad de configurar la IP manualmente. El segundo es verificar la cobertura Wi-Fi y el estado de los puntos de acceso en las aulas y ambientes reportados. El tercero es verificar la configuración DHCP, que asigna direcciones IP automáticamente a los dispositivos \cite{rfc2131}, y las concesiones de direcciones en los laboratorios. El cuarto es medir la velocidad, la latencia y la pérdida de tramas en horas punta y valle, con una metodología de pruebas de rendimiento de dispositivos de red \cite{rfc2544}, y relacionar esas mediciones con la cantidad de usuarios concurrentes. El quinto es comparar las alternativas de solución, incluida la optimización de la red existente, con criterios comunes antes de definir inversiones.

% SUB-SECCIÓN
\subsection{Relación entre problema y objetivos}

	La falta de señal o el acceso no utilizable en algunas aulas se atiende con el registro por aula y con la prueba de cobertura y del estado de los puntos de acceso. La necesidad de configurar la IP manualmente en los laboratorios se atiende con pruebas con varios dispositivos y con la revisión de DHCP. Los cortes de conexión se atienden con el registro de episodios y su duración, y con la revisión de los registros de enlace y autenticación. La lentitud en la red institucional se atiende con pruebas de velocidad y latencia en horas punta y valle. Las demoras en las tareas académicas se atienden con el registro de las actividades afectadas y con la comparación de alternativas.

% SUB-SECCIÓN
\subsection{Indicadores verificables}

	Se proponen tres indicadores. El primero es la cobertura Wi-Fi verificada, es decir, el porcentaje de ambientes con señal utilizable en las pruebas de campo. El segundo es la cantidad de episodios de corte por semana en el pabellón y su duración promedio, según el registro. El tercero es el tiempo que transcurre desde la conexión hasta el acceso a Internet en los laboratorios, medido en varios dispositivos. Las metas de los tres indicadores se definirán después del diagnóstico, porque todavía no existe una línea base.

	Los objetivos no buscan únicamente comprar equipos ni desarrollar una aplicación. Su resultado esperado es una mejora medible del acceso a Internet, y la solución se elegirá después de verificar las causas.
